\documentclass[journal]{IEEEtran}
\usepackage{amsmath,amssymb}
\usepackage{booktabs}
\usepackage{graphicx}
\usepackage[hidelinks]{hyperref}

\newcommand{\agecal}{\textit{agentic-eCAL}}
\newcommand{\ecal}{\textit{eCAL}}

\begin{document}

\title{Adaptive Model Selection in Agentic Workflows:\\An Energy Accounting}
\author{\IEEEauthorblockN{Anonymous Authors}\thanks{Working draft, not for circulation.}}
\maketitle

\begin{abstract}
Agentic workflows execute a graph of language-model calls whose steps differ markedly in the
reasoning they demand: a planner emits a short task decomposition, whereas a writer must
synthesise a long answer over an accumulated transcript. Serving every step with the same model
therefore over-provisions most of them. This note prices adaptive model selection with the
two-rate call model of the companion paper, and reports the energy available from assigning each
step the smallest model that could plausibly serve it. It deliberately stops short of claiming a
quality-neutral saving, because the measurements available to us characterise energy but not task
success; what follows is an upper bound on the opportunity, together with the experiment that
would be required to convert it into a result.
\end{abstract}

\section{Model}
Each call is priced by the two-rate model, with a compute-bound prefill rate and a
memory-bound decode rate that amortises the weight read over the serving batch $b$:
\begin{equation}
E_{\mathrm{call}}(p_{\mathrm{in}},p_{\mathrm{out}};b)
 = c_{\mathrm{pre}}\,p_{\mathrm{in}} + c_{\mathrm{dec}}(b)\,p_{\mathrm{out}} .
\label{eq:tworate}
\end{equation}
Both rates scale with the active parameter count $N$: $c_{\mathrm{pre}}=2NP/(\eta_{\mathrm{pre}}\Pi)$
and the weight term of $c_{\mathrm{dec}}$ as $N\beta/b$. Model selection is therefore a direct
energy lever, and its magnitude follows from the parameter ratio rather than from any property of
the workflow.

\begin{table}[t]
\caption{Two-rate coefficients on an A100 at $b=64$, derived from datasheet quantities.}
\label{tab:coeff}
\centering
\small
\begin{tabular}{@{}lrrr@{}}
\toprule
Model & $N$ [B] & $c_{\mathrm{pre}}$ [J/tok] & $c_{\mathrm{dec}}(64)$ [J/tok] \\
\midrule
Qwen2.5-7B  &  7.6 & 0.0230 & 0.0828 \\
Llama-3.1-8B &  8.0 & 0.0242 & 0.1071 \\
Llama-3 70B & 70.6 & 0.2130 & 0.7101 \\
\bottomrule
\end{tabular}
\end{table}

\section{Opportunity}
Table~\ref{tab:assign} prices the four-agent retrieval-augmented workflow of the companion paper
over $K=3$ rounds at $b=64$, under homogeneous and mixed assignments. Serving every step with the
70B model costs $7.6\times$ the all-8B assignment. Reserving the large model for the writer alone
recovers most of that: $1808$\,J against $4971$\,J, a $2.7\times$ reduction, while leaving the
single step that produces the user-visible output unchanged.

\begin{table}[t]
\caption{Workflow energy under model assignment. Four-agent RAG workflow, $K=3$, $b=64$, A100.}
\label{tab:assign}
\centering
\small
\begin{tabular}{@{}lrr@{}}
\toprule
Assignment & $E_W$ [J] & vs.\ all-70B \\
\midrule
All steps on Llama-3 70B                     & 4971 & $1.00\times$ \\
70B for the analyst and writer               & 2765 & $0.56\times$ \\
70B for the writer only                      & 1808 & $0.36\times$ \\
All steps on Llama-3.1 8B                    &  653 & $0.13\times$ \\
\bottomrule
\end{tabular}
\end{table}

\section{What this does not establish}
The assignments above are priced, not validated. Nothing in the measurements underlying
Table~\ref{tab:coeff} indicates whether an 8B model planning a task yields a workflow whose output
matches one planned by a 70B model, and the energy ordering is therefore an upper bound on a
saving whose feasibility is untested. Two properties of agentic workflows make the question harder
than it is for a single inference. Errors propagate: a weak plan is inherited by every later step,
so per-step quality is not the relevant measure. And the steps are not interchangeable: the
literature on routing between models addresses one-shot queries, whereas here the cost of a poor
choice appears several steps downstream.

Converting this note into a result requires an experiment we have not run: the same workflow
executed under each assignment on a task with a machine-checkable outcome, with success rate
recorded alongside energy, so that the axis reported is joules per \emph{successful} workflow
rather than joules per workflow. Until then the honest claim is the one made here---that model
assignment is the largest single energy lever available to an agentic deployment, worth up to
$7.6\times$, and that whether it can be exercised without loss is an open question.

\end{document}
